% preamble-palabras.tex
% -----------------------------------------------------------------------
% Lo propio del cuaderno de primeras palabras (palabras.tex,
% palabras-bw.tex), cargado DESPUÉS de preamble.tex: la letra, la paleta,
% las cajas base, \cabeceraDia, \casilla y la actividad "traza" son las
% mismas que en el cuaderno de frases -- aquí solo lo que cambia en este
% nivel:
%
%   - la página: la caja de actividad llena todo lo que queda de página
%     (tcolorbox `height fill`), porque en este nivel casi todas las
%     actividades terminan dibujando, y el sitio para dibujar es todo el
%     que haya -- nunca una caja pequeña con media página en blanco
%     debajo (ver notes/02-revision-y-plan.md, punto 5). Para eso el pie
%     "Día N / 260" tiene que salir de la página de texto y vivir en el
%     pie de página de verdad (estilo `piepalabras`), si no la caja lo
%     empujaría a una segunda página;
%   - la caja de lectura: tarjetas de palabra (sílabas arriba, palabra
%     entera debajo), una, dos o tres por página según el trimestre;
%   - las actividades nuevas: escribe, encuentra, palmadas, une, sí/no.
% -----------------------------------------------------------------------

\usepackage{pdfrender}

% El pie de página de verdad necesita sitio por debajo del texto: 22 mm
% de margen inferior con el pie a 10 mm del texto deja la línea del pie
% a ~12 mm del borde del papel, dentro de lo que imprime cualquier
% impresora doméstica.
\geometry{bottom=22mm, footskip=10mm}

\makeatletter
\def\ps@piepalabras{%
  \let\@oddhead\@empty\let\@evenhead\@empty
  \def\@oddfoot{\hfil\footnotesize\color{colorGris}\lblDia~\numeroDia{\palabrapagina@actual}~/ \totaldias\hfil}%
  \let\@evenfoot\@oddfoot
}

% La página de un día, en este cuaderno. Misma cabecera que el de
% frases (\cabeceraDia, preamble.tex) y el mismo \label{dia:N}, así que
% tools/check_pages.py funciona igual sobre los dos .aux. \gdef y no
% \def: el pie se compone en la rutina de salida, y aunque esa rutina
% corre todavía dentro del grupo del entorno, así no depende de ello.
\newenvironment{palabrapagina}[4]{%
  \gdef\palabrapagina@actual{#1}%
  \thispagestyle{piepalabras}%
  \cabeceraDia{#1}{#2}{#3}{#4}%
}{%
  \par\newpage
}
\makeatother

% =========================================================================
% Caja de lectura: tarjetas de palabra
% =========================================================================
% Tamaño de la palabra y de sus sílabas, según el trimestre. T1: una
% palabra enorme; T3: tres por página, más pequeñas para que quepan en
% una fila. T4 no tiene tarjetas -- ya es una frase, sin sílabas (ver
% \fraseDia).
\newcommand{\fuentePalabra}[1]{%
  \ifcase#1\relax
  \or \fontsize{64}{70}\selectfont
  \or \fontsize{46}{52}\selectfont
  \or \fontsize{31}{37}\selectfont
  \fi
}
\newcommand{\fuenteSilabas}[1]{%
  \ifcase#1\relax
  \or \fontsize{28}{32}\selectfont
  \or \fontsize{24}{28}\selectfont
  \or \fontsize{18}{22}\selectfont
  \fi
}
% Separador entre sílabas: un punto centrado con aire a los dos lados,
% más fácil de ver para un niño que un guion.
\newcommand{\sep}{\hspace{0.28em}\textperiodcentered\hspace{0.28em}}

% Una tarjeta: primero las sílabas (en color, más pequeñas), después la
% palabra entera -- el orden en que se leen: sílaba a sílaba, luego de
% corrido.
%
% Una palabra larga de T3 (va·ca·cio·nes, ra·ton·ci·to) no cabe a ese
% tamaño en una columna de un tercio de página: \ajustado la encoge
% justo lo necesario para que quepa, en vez de dejar que se salga de la
% caja. Las palabras que ya caben no se tocan.
\newsavebox{\cajaAjuste}
\newcommand{\ajustado}[1]{%
  \sbox{\cajaAjuste}{#1}%
  \ifdim\wd\cajaAjuste>\linewidth
    \resizebox{\linewidth}{!}{\usebox{\cajaAjuste}}%
  \else
    \usebox{\cajaAjuste}%
  \fi
}
\newcommand{\tarjeta}[3]{%
  {\fuenteSilabas{#1}\color{colorLectura}\ajustado{#2}\par}%
  \vspace{1.5mm}%
  {\fuentePalabra{#1}\ajustado{#3}\par}%
}

% Un nombre que se lee de un golpe (Toby en «First Words», Lucía en
% «Pierwsze słowa»): la palabra entera, sin partir en sílabas ni botones
% de sonidos, con el mismo aire debajo que una tarjeta con botones.
\newcommand{\tarjetaEntera}[2]{%
  {\fuentePalabra{#1}\ajustado{#2}\par\vspace{0.3em}}%
}

% La frase del día en T4: una sola línea, grande, sin sílabas. Una
% frase de cuatro palabras largas ("Lápices y cuadernos nuevos.") se
% encoge un poco antes que partirse en dos líneas: a este nivel, una
% frase que salta de línea a mitad se lee como dos.
\newcommand{\fraseDia}[1]{{\centering\fontsize{36}{44}\selectfont\ajustado{#1}\par}}

% El viernes: las palabras (o frases) de toda la semana, cada una con
% una casilla para marcar si ya la lee sola. Tamaño según cuántas son.
\newcommand{\fuenteSemana}[1]{%
  \ifcase#1\relax
  \or \fontsize{30}{36}\selectfont
  \or \fontsize{25}{31}\selectfont
  \or \fontsize{21}{27}\selectfont
  \or \fontsize{22}{30}\selectfont
  \fi
}
\newcommand{\casillaGrande}{\raisebox{0.05ex}{\framebox[5mm][c]{\rule{0pt}{4mm}}}}
\newcommand{\palabraSemana}[1]{\casillaGrande\hspace{2mm}#1}

% =========================================================================
% Cajas de actividad: las mismas que en el cuaderno de frases (colores,
% títulos), pero todas con `height fill` -- llenan lo que queda de
% página. El \vspace{5mm} que las separa de la caja de lectura lo pone
% la plantilla del día (tools/gen_palabras.py).
% =========================================================================
% Una caja de altura fija no respeta un \vfill dentro (tcolorbox mide
% el contenido a su altura natural antes de colocarlo), así que lo que
% tiene que ir abajo del todo (la línea de escribir de "adivina", el
% cartel de "repasa") o centrado en el hueco (el dibujo de "completa")
% va en la parte de abajo de la caja (\tcblower), con `abajo` o
% `centrado abajo`: todo el sitio sobrante se lo lleva la otra parte.
\tcbset{
  cajapalabras/.style={cajabase, height fill},
  abajo/.style={space to upper, segmentation hidden, middle=2mm},
  centrado abajo/.style={space to lower, valign lower=center, segmentation hidden, middle=2mm},
}
\newtcolorbox{cajaPalDibuja}[1][]{cajapalabras, #1,
  colback=colorDibujaClaro, colframe=colorDibuja, colbacktitle=colorDibuja,
  title=\lblDibuja}
\newtcolorbox{cajaPalCompleta}[1][]{cajapalabras, #1,
  colback=colorCompletaClaro, colframe=colorCompleta, colbacktitle=colorCompleta,
  title=\lblCompleta}
\newtcolorbox{cajaPalTraza}[1][]{cajapalabras, #1,
  colback=colorRespondeClaro, colframe=colorResponde, colbacktitle=colorResponde,
  title=\lblTraza}
\newtcolorbox{cajaPalEscribe}[1][]{cajapalabras, #1,
  colback=colorRespondeClaro, colframe=colorResponde, colbacktitle=colorResponde,
  title=\lblEscribe}
\newtcolorbox{cajaPalEncuentra}[1][]{cajapalabras, #1,
  colback=colorCompletaClaro, colframe=colorCompleta, colbacktitle=colorCompleta,
  title=\lblEncuentra}
\newtcolorbox{cajaPalPalmadas}[1][]{cajapalabras, #1,
  colback=colorCompletaClaro, colframe=colorCompleta, colbacktitle=colorCompleta,
  title=\lblPalmadas}
\newtcolorbox{cajaPalAdivina}[1][]{cajapalabras, #1,
  colback=colorRelacionaClaro, colframe=colorRelaciona, colbacktitle=colorRelaciona,
  title=\lblAdivina}
\newtcolorbox{cajaPalUne}[1][]{cajapalabras, #1,
  colback=colorRelacionaClaro, colframe=colorRelaciona, colbacktitle=colorRelaciona,
  title=\lblUne}
\newtcolorbox{cajaPalSiNo}[1][]{cajapalabras, #1,
  colback=colorRespondeClaro, colframe=colorResponde, colbacktitle=colorResponde,
  title=\lblSiNo}
\newtcolorbox{cajaPalRelee}[1][]{cajapalabras, #1,
  colback=colorCreaClaro, colframe=colorCrea, colbacktitle=colorCrea,
  title=\lblRelee}
\newtcolorbox{cajaPalRepasa}[1][]{cajapalabras, #1,
  colback=colorCreaClaro, colframe=colorCrea, colbacktitle=colorCrea,
  title=\lblRepasa}

% Una instrucción para el adulto: pequeña y gris, como en el cuaderno de
% frases ("Une cada nombre..."). Lo que lee la niña o el niño va aparte,
% siempre grande.
\newcommand{\instruccion}[1]{{\footnotesize\color{colorGris}#1\par}\vspace{3mm}}

% El remate de dibujo que llevan casi todas las actividades: una línea
% de instrucción y, debajo, todo el sitio que quede en la caja.
\newcommand{\ahoraDibuja}[1]{\par\vspace{6mm}{\footnotesize\color{colorGris}\lblPalAhoraDibuja}\ {\large #1}\par}

% Una línea para escribir, a lo ancho de la caja, con aire encima.
\newcommand{\lineaEscribir}{\par\vspace{16mm}\noindent\rule{\linewidth}{0.5pt}\par}

% Palabra hueca, para repasarla con el lápiz por dentro (actividad
% "escribe"): solo el contorno de cada letra, en gris -- la misma letra
% Andika del resto del cuaderno, no una fuente punteada aparte.
\newcommand{\palabraHueca}[1]{%
  {\fontsize{76}{96}\selectfont
   \textpdfrender{TextRenderingMode=Stroke, LineWidth=0.9pt, StrokeColor=colorGris}{#1}\par}%
}

% Un círculo vacío para colorear (actividad "palmadas").
\newcommand{\circuloPalmada}{\tikz[baseline=-0.6ex]\draw[line width=1pt, colorGris] (0,0) circle (4.2mm);}
